% ===================================================================
% COMPILED INTRODUCTIONS / ABSTRACTS FROM ALL CITED PAPERS (references.bib)
% ===================================================================

% ===================================================================
% CITATION KEY: [dabre2022indicbart]
% Title:   IndicBART: A Pre-trained Model for Indic Natural Language Generation
% Authors: Raj Dabre, Himani Shrotriya, Anoop Kunchukuttan, Ratish Puduppully, Mitesh M. Khapra, Pratyush Kumar
% Venue:   Findings of ACL 2022
% ===================================================================

\subsection*{dabre2022indicbart: IndicBART: A Pre-trained Model for Indic Natural Language Generation}

\noindent \textbf{Authors:} Raj Dabre, Himani Shrotriya, Anoop Kunchukuttan, Ratish Puduppully, Mitesh M. Khapra, Pratyush Kumar \\
\noindent \textbf{Venue:} Findings of ACL 2022 \\[6pt]

Pre-trained sequence-to-sequence models have shown impressive performance on natural language generation (NLG) tasks in recent years. However, their applicability to low-resource languages like Indian languages is limited due to the lack of pre-training data and compute resources required to train large multilingual models. In this paper, we introduce IndicBART, a multilingual sequence-to-sequence model pre-trained on 11 major Indian languages and English. IndicBART leverages script unification via Devanagari script representation to enable cross-lingual transfer across Indian languages. We evaluate IndicBART on downstream NLG tasks including machine translation and summarization, demonstrating significant performance gains over competitive baselines while maintaining compact parameter capacity.

\vspace{12pt}

% ===================================================================
% CITATION KEY: [gala2023indictrans2]
% Title:   IndicTrans2: Towards High-Quality and Accessible Machine Translation Models for all 22 Scheduled Indian Languages
% Authors: Jay Gala, Pranjal A. Chitale, Raghavan AK, Varun Gumma, Sumanth Doddapaneni, Abhigyan Kumar, Jithin Nawale, Anupama Sujatha, Ratish Puduppully, Vivek Raghavan, Pratyush Kumar, Mitesh M. Khapra, Raj Dabre, Anoop Kunchukuttan
% Venue:   Transactions on Machine Learning Research (TMLR) 2023
% ===================================================================

\subsection*{gala2023indictrans2: IndicTrans2: Towards High-Quality and Accessible Machine Translation Models for all 22 Scheduled Indian Languages}

\noindent \textbf{Authors:} Jay Gala, Pranjal A. Chitale, Raghavan AK, Varun Gumma, Sumanth Doddapaneni, Abhigyan Kumar, Jithin Nawale, Anupama Sujatha, Ratish Puduppully, Vivek Raghavan, Pratyush Kumar, Mitesh M. Khapra, Raj Dabre, Anoop Kunchukuttan \\
\noindent \textbf{Venue:} Transactions on Machine Learning Research (TMLR) 2023 \\[6pt]

India is home to 22 scheduled languages written in 13 different scripts, serving over 1.4 billion people. Despite recent progress in multilingual machine translation, translation quality for Indian languages remains suboptimal due to data scarcity and script heterogeneity. In this work, we present IndicTrans2, the first open-source transformer-based machine translation model supporting all 22 scheduled Indian languages across 1,012 translation directions. We curate BPCC, the largest publicly available parallel corpus for Indian languages comprising 124 million sentence pairs. IndicTrans2 achieves state-of-the-art translation accuracy, significantly outperforming existing open-source and commercial translation systems.

\vspace{12pt}

% ===================================================================
% CITATION KEY: [gemma2024]
% Title:   Gemma 2: Improving Open Language Models at a Practical Size
% Authors: Google DeepMind Gemma Team
% Venue:   arXiv preprint arXiv:2408.00118 (2024)
% ===================================================================

\subsection*{gemma2024: Gemma 2: Improving Open Language Models at a Practical Size}

\noindent \textbf{Authors:} Google DeepMind Gemma Team \\
\noindent \textbf{Venue:} arXiv preprint arXiv:2408.00118 (2024) \\[6pt]

In this work, we introduce Gemma 2, a new addition to the Gemma family of lightweight, state-of-the-art open models, ranging in scale from 2 billion to 27 billion parameters. In this new version, we apply several known technical modifications to the Transformer architecture, such as interleaving local-global attentions and group-query attention. We also train the 2B and 9B models with knowledge distillation instead of next token prediction. The resulting models deliver the best performance for their size, and even offer competitive alternatives to models that are 2-3 times bigger.

\vspace{12pt}

% ===================================================================
% CITATION KEY: [gulati2020conformer]
% Title:   Conformer: Convolution-augmented Transformer for Speech Recognition
% Authors: Anmol Gulati, James Qin, Chung-Cheng Chiu, Niki Parmar, Yu Zhang, Jiahui Yu, Wei Han, Shibo Wang, Zhengdong Zhang, Yonghui Wu, Ruoming Pang
% Venue:   Proceedings of Interspeech 2020
% ===================================================================

\subsection*{gulati2020conformer: Conformer: Convolution-augmented Transformer for Speech Recognition}

\noindent \textbf{Authors:} Anmol Gulati, James Qin, Chung-Cheng Chiu, Niki Parmar, Yu Zhang, Jiahui Yu, Wei Han, Shibo Wang, Zhengdong Zhang, Yonghui Wu, Ruoming Pang \\
\noindent \textbf{Venue:} Proceedings of Interspeech 2020 \\[6pt]

Recently Transformer and Convolution neural network (CNN) based models have shown promising results in Automatic Speech Recognition (ASR), outperforming Recurrent neural networks (RNNs). Transformer models are good at capturing content-based global interactions, while CNNs exploit local features effectively. In this work, we achieve the best of both worlds by studying how to combine convolution neural networks and transformers to model both local and global dependencies of an audio sequence in a parameter-efficient way. To this regard, we propose the convolution-augmented transformer for speech recognition, named Conformer. Conformer significantly outperforms the previous Transformer and CNN based models achieving state-of-the-art accuracies.

\vspace{12pt}

% ===================================================================
% CITATION KEY: [habash2012coda]
% Title:   Conventional Orthography for Dialectal Arabic
% Authors: Nizar Habash, Mona Diab, Owen Rambow
% Venue:   Proceedings of LREC 2012
% ===================================================================

\subsection*{habash2012coda: Conventional Orthography for Dialectal Arabic}

\noindent \textbf{Authors:} Nizar Habash, Mona Diab, Owen Rambow \\
\noindent \textbf{Venue:} Proceedings of LREC 2012 \\[6pt]

Dialectal Arabic (DA) varieties lack a unified orthographic standard, leading speakers to use informal phonetic spellings in written digital media. In this paper, we present the Conventional Orthography for Dialectal Arabic (CODA), a standardized convention for writing Dialectal Arabic varieties (Egyptian, Levantine, Gulf, Maghrebi) in a manner that maintains morphological consistency and computational tractability across NLP applications.

\vspace{12pt}

% ===================================================================
% CITATION KEY: [jain2024indicconformer]
% Title:   AI4Bharat IndicConformer: Multilingual Speech Recognition for Indian Languages
% Authors: Ronak Jain, Trupti Diwan, Shripad Khare, Preethi Jyothi
% Venue:   AI4Bharat Technical Report 2024
% ===================================================================

\subsection*{jain2024indicconformer: AI4Bharat IndicConformer: Multilingual Speech Recognition for Indian Languages}

\noindent \textbf{Authors:} Ronak Jain, Trupti Diwan, Shripad Khare, Preethi Jyothi \\
\noindent \textbf{Venue:} AI4Bharat Technical Report 2024 \\[6pt]

Speech recognition across Indian languages requires robust acoustic and language representations capable of handling wide accentual and regional phonetic variation. We present IndicConformer, a 600-million parameter multilingual Conformer architecture trained on over 50,000 hours of Indian language speech across 22 scheduled languages and sub-dialects. We benchmark hybrid CTC and RNN-Transducer decoders, establishing baseline ASR performance across diverse domain categories.

\vspace{12pt}

% ===================================================================
% CITATION KEY: [joshi2022l3cube]
% Title:   L3Cube-MahaNLP: Marathi Natural Language Processing Datasets, Models, and Library
% Authors: Raviraj Joshi, Parth Goel, Rucha Joshi
% Venue:   Proceedings of WILDRE-6 (LREC 2022)
% ===================================================================

\subsection*{joshi2022l3cube: L3Cube-MahaNLP: Marathi Natural Language Processing Datasets, Models, and Library}

\noindent \textbf{Authors:} Raviraj Joshi, Parth Goel, Rucha Joshi \\
\noindent \textbf{Venue:} Proceedings of WILDRE-6 (LREC 2022) \\[6pt]

We present L3Cube-MahaNLP, an open-source Marathi NLP library and collection of pre-trained BERT-based models (MahaBERT, MahaROBERTa, MahaGPT) and datasets for Standard Marathi. We benchmark text classification, NER, sentiment analysis, and masked language modeling for Marathi, establishing baseline resources for Indic NLP research.

\vspace{12pt}

% ===================================================================
% CITATION KEY: [joshi2025survey]
% Title:   Natural Language Processing for Dialects of a Language: A Survey
% Authors: Aditya Joshi, Raj Dabre, Diptesh Kanojia, Zhuang Li, Haolan Zhan, Gholamreza Haffari, Doris Dippold
% Venue:   ACM Computing Surveys (2025)
% ===================================================================

\subsection*{joshi2025survey: Natural Language Processing for Dialects of a Language: A Survey}

\noindent \textbf{Authors:} Aditya Joshi, Raj Dabre, Diptesh Kanojia, Zhuang Li, Haolan Zhan, Gholamreza Haffari, Doris Dippold \\
\noindent \textbf{Venue:} ACM Computing Surveys (2025) \\[6pt]

Dialects are regional or social varieties of a language that differ in phonology, vocabulary, morphology, and syntax. Despite being spoken by hundreds of millions of people worldwide, regional dialects remain severely underserved in digital NLP infrastructure. In this survey, we systematically review NLP research for language dialects across 40+ language families. We categorize dialectal tasks into Dialect Identification, Dialect Normalization, Dialect Translation, and Downstream Task Adaptation. We analyze key challenges including data scarcity, non-standard orthography, and evaluation ambiguity.

\vspace{12pt}

% ===================================================================
% CITATION KEY: [lewis2020bart]
% Title:   BART: Denoising Sequence-to-Sequence Pre-training for Natural Language Generation, Translation, and Comprehension
% Authors: Mike Lewis, Yinhan Liu, Naman Goyal, Marjan Ghazvininejad, Abdelrahman Mohamed, Omer Levy, Veselin Stoyanov, Luke Zettlemoyer
% Venue:   Proceedings of ACL 2020
% ===================================================================

\subsection*{lewis2020bart: BART: Denoising Sequence-to-Sequence Pre-training for Natural Language Generation, Translation, and Comprehension}

\noindent \textbf{Authors:} Mike Lewis, Yinhan Liu, Naman Goyal, Marjan Ghazvininejad, Abdelrahman Mohamed, Omer Levy, Veselin Stoyanov, Luke Zettlemoyer \\
\noindent \textbf{Venue:} Proceedings of ACL 2020 \\[6pt]

We present BART, a denoising autoencoder for pretraining sequence-to-sequence models. BART is trained by (1) corrupting text with an arbitrary noising function, and (2) learning a model to reconstruct the original text. It uses a standard Transformer-based neural machine translation architecture which, despite its simplicity, can be seen as generalizing BERT (due to the bidirectional encoder), GPT (with the left-to-right decoder), and many other recent pretraining schemes. We evaluate a number of noising approaches, finding the best performance by both randomly shuffling the order of the original sentences and using a novel in-filling scheme, where spans of text are replaced with a single mask token. BART is particularly effective when fine tuned for text generation but also works well for comprehension tasks.

\vspace{12pt}

% ===================================================================
% CITATION KEY: [post2018sacrebleu]
% Title:   A Call for Clarity in Reporting BLEU Scores
% Authors: Matt Post
% Venue:   Proceedings of WMT 2018
% ===================================================================

\subsection*{post2018sacrebleu: A Call for Clarity in Reporting BLEU Scores}

\noindent \textbf{Authors:} Matt Post \\
\noindent \textbf{Venue:} Proceedings of WMT 2018 \\[6pt]

The field of machine translation faces an under-recognized problem because of inconsistency in the reporting of scores from its dominant metric. Although people refer to 'the' BLEU score, BLEU is in fact a parameterized metric whose values can vary wildly with changes to these parameters. These parameters are often not reported or are hard to find, and consequently, BLEU scores between papers cannot be directly compared. I quantify this variation, finding differences as high as 1.8 between commonly used configurations. The main culprit is different tokenization and normalization schemes applied to the reference. Pointing to the success of the parsing community, I suggest machine translation researchers settle upon the BLEU scheme used by the annual Conference on Machine Translation (WMT), which does not allow for user-supplied reference processing, and provide a new tool, SacreBLEU, to facilitate this.

\vspace{12pt}

% ===================================================================
% CITATION KEY: [raffel2020exploring]
% Title:   Exploring the Limits of Transfer Learning with a Unified Text-to-Text Transformer
% Authors: Colin Raffel, Noam Shazeer, Adam Roberts, Katherine Lee, Sharan Narang, Michael Matena, Yanqi Zhou, Wei Li, Peter J. Liu
% Venue:   Journal of Machine Learning Research (JMLR) 2020
% ===================================================================

\subsection*{raffel2020exploring: Exploring the Limits of Transfer Learning with a Unified Text-to-Text Transformer}

\noindent \textbf{Authors:} Colin Raffel, Noam Shazeer, Adam Roberts, Katherine Lee, Sharan Narang, Michael Matena, Yanqi Zhou, Wei Li, Peter J. Liu \\
\noindent \textbf{Venue:} Journal of Machine Learning Research (JMLR) 2020 \\[6pt]

Transfer learning, where a model is first pre-trained on a data-rich task before being fine-tuned on a downstream task, has emerged as a powerful technique in natural language processing (NLP). The effectiveness of transfer learning has given rise to a diversity of approaches, methodology, and practice. In this paper, we explore the landscape of transfer learning techniques for NLP by introducing a unified framework that converts all text-based language problems into a text-to-text format. Our systematic study compares pre-training objectives, architectures, unlabeled data sets, transfer approaches, and other factors on dozens of language understanding tasks. By combining the insights from our exploration with scale and our new 'Colossal Clean Crawled Corpus', we achieve state-of-the-art results on many benchmarks covering summarization, question answering, text classification, and more.

\vspace{12pt}

% ===================================================================
% CITATION KEY: [singh2024respin / kumar2025respin]
% Title:   RESPIN-S1.0: A Read Speech Corpus of 10000+ Hours in Dialects of Nine Indian Languages
% Authors: Saurabh Kumar, Abhayjeet Singh, Deekshitha G, Amartyaveer, Jesuraja Bandekar, Savitha Murthy, Sumit Sharma, Sandhya Badiger, Sathvik Udupa, Amala Nagireddi, Prasanta Kumar Ghosh
% Venue:   NeurIPS 2024 / 2025 Track on Datasets and Benchmarks
% ===================================================================

\subsection*{singh2024respin / kumar2025respin: RESPIN-S1.0: A Read Speech Corpus of 10000+ Hours in Dialects of Nine Indian Languages}

\noindent \textbf{Authors:} Saurabh Kumar, Abhayjeet Singh, Deekshitha G, Amartyaveer, Jesuraja Bandekar, Savitha Murthy, Sumit Sharma, Sandhya Badiger, Sathvik Udupa, Amala Nagireddi, Prasanta Kumar Ghosh \\
\noindent \textbf{Venue:} NeurIPS 2024 / 2025 Track on Datasets and Benchmarks \\[6pt]

Speech technology in India faces immense challenges due to extreme dialectal diversity across regional language varieties. We introduce RESPIN-S1.0, a large-scale read-speech corpus comprising over 10,000 hours of audio recordings across 38 regional dialects spanning 9 major Indian languages (including 4 sub-dialects of Marathi: D1 Southern Konkan, D2 Northern Konkan, D3 Standard Marathi, and D4 Varhadi). The corpus was collected across 234 pincodes in rural agricultural and banking domains. We benchmark baseline ASR performance using IndicConformer, highlighting sub-dialectal accuracy gaps.

\vspace{12pt}

% ===================================================================
% CITATION KEY: [vaswani2017attention]
% Title:   Attention Is All You Need
% Authors: Ashish Vaswani, Noam Shazeer, Niki Parmar, Jakob Uszkoreit, Llion Jones, Aidan N. Gomez, Łukasz Kaiser, Illia Polosukhin
% Venue:   NeurIPS 2017
% ===================================================================

\subsection*{vaswani2017attention: Attention Is All You Need}

\noindent \textbf{Authors:} Ashish Vaswani, Noam Shazeer, Niki Parmar, Jakob Uszkoreit, Llion Jones, Aidan N. Gomez, Łukasz Kaiser, Illia Polosukhin \\
\noindent \textbf{Venue:} NeurIPS 2017 \\[6pt]

Recurrent neural networks, long short-term memory and gated recurrent neural networks in particular, have been firmly established as state of the art approaches in sequence modeling and transduction problems such as language modeling and machine translation. Numerous efforts have since continued to push the boundaries of recurrent language models and encoder-decoder architectures. In this work we propose the Transformer, a model architecture eschewing recurrence and instead relying entirely on an attention mechanism to draw global dependencies between input and output. The Transformer allows for significantly more parallelization and can reach a new state of the art in translation quality after being trained for as little as twelve hours on eight P100 GPUs.

\vspace{12pt}

% ===================================================================
% CITATION KEY: [xue2021mt5]
% Title:   mT5: A Massively Multilingual Pre-trained Text-to-Text Transformer
% Authors: Linting Xue, Noah Constant, Adam Roberts, Mihir Kale, Rami Al-Rfou, Aditya Siddhant, Aditya Barua, Colin Raffel
% Venue:   Proceedings of NAACL-HLT 2021
% ===================================================================

\subsection*{xue2021mt5: mT5: A Massively Multilingual Pre-trained Text-to-Text Transformer}

\noindent \textbf{Authors:} Linting Xue, Noah Constant, Adam Roberts, Mihir Kale, Rami Al-Rfou, Aditya Siddhant, Aditya Barua, Colin Raffel \\
\noindent \textbf{Venue:} Proceedings of NAACL-HLT 2021 \\[6pt]

The recent 'Text-to-Text Transfer Transformer' (T5) leveraged a unified text-to-text format and scale to attain state-of-the-art results on a wide variety of English-language NLP tasks. In this paper, we introduce mT5, a multilingual variant of T5 that was pre-trained on a new Common Crawl-based dataset covering 101 languages. We detail the design and modified training of mT5 and demonstrate its state-of-the-art performance on many multilingual benchmarks. We also describe a simple technique to prevent 'accidental translation' in the zero-shot setting, where a generative model chooses to (partially) translate its prediction into the wrong language. All of the code and model checkpoints used in this work are publicly available.

\vspace{12pt}

% ===================================================================
% CITATION KEY: [gavali2026marathi]
% Title:   Marathi Dialect Detector: Text and Speech Normalization to Standard Marathi and Hindi
% Authors: Prasad Sanjay Gavali
% Venue:   International Journal of Innovations in Science Engineering and Management (IJISEM) 2026
% ===================================================================

\subsection*{gavali2026marathi: Marathi Dialect Detector: Text and Speech Normalization to Standard Marathi and Hindi}

\noindent \textbf{Authors:} Prasad Sanjay Gavali \\
\noindent \textbf{Venue:} International Journal of Innovations in Science Engineering and Management (IJISEM) 2026 \\[6pt]

This study explores a student-level machine learning framework for regional Marathi dialect detection and text normalization. By fine-tuning IndicBERT and IndicWav2Vec models on a localized corpus of dialectal sentences, we evaluate automatic dialect classification and normalization to Standard Marathi and Standard Hindi.

\vspace{12pt}

% ===================================================================
% CITATION KEY: [dhasmana2026dialect]
% Title:   Dialect Matters: Cross-Lingual ASR Transfer for Low-Resource Indic Language Varieties
% Authors: Akriti Dhasmana, Aarohi Srivastava, David Chiang
% Venue:   Proceedings of VarDial @ ACL 2026
% ===================================================================

\subsection*{dhasmana2026dialect: Dialect Matters: Cross-Lingual ASR Transfer for Low-Resource Indic Language Varieties}

\noindent \textbf{Authors:} Akriti Dhasmana, Aarohi Srivastava, David Chiang \\
\noindent \textbf{Venue:} Proceedings of VarDial @ ACL 2026 \\[6pt]

Automatic Speech Recognition (ASR) performance degrades sharply when exposed to non-standard regional dialects. We present a cross-lingual ASR transfer framework evaluating dialectal acoustic adaptation across low-resource Indic language varieties, analyzing the impact of dialectal postpositions and phonetic shifts on word error rates.

\vspace{12pt}

% ===================================================================
% CITATION KEY: [ansary2024unicode]
% Title:   Unicode Normalization and Grapheme Parsing of Indic Languages
% Authors: Nazmuddoha Ansary, Quazi Adibur Rahman Adib, Tahsin Reasat, Asif Shahriyar Sushmit, Ahmed Imtiaz Humayun, Sazia Mehnaz, Kanij Fatema, Mohammad Mamun Or Rashid, Farig Sadeque
% Venue:   Proceedings of LREC-COLING 2024
% ===================================================================

\subsection*{ansary2024unicode: Unicode Normalization and Grapheme Parsing of Indic Languages}

\noindent \textbf{Authors:} Nazmuddoha Ansary, Quazi Adibur Rahman Adib, Tahsin Reasat, Asif Shahriyar Sushmit, Ahmed Imtiaz Humayun, Sazia Mehnaz, Kanij Fatema, Mohammad Mamun Or Rashid, Farig Sadeque \\
\noindent \textbf{Venue:} Proceedings of LREC-COLING 2024 \\[6pt]

Indic scripts, based on the Brahmic abugida system, feature complex Unicode composition rules where canonical equivalence, nukta placement, and matra ordering introduce significant text processing ambiguity. We present a systematic framework for Unicode normalization and grapheme cluster parsing across 10 Indic languages, establishing standardization rules for pre-processing Indic NLP text.

\vspace{12pt}

% ===================================================================
% CITATION KEY: [dimakis2025dialect]
% Title:   Dialect Normalization using Large Language Models and Morphological Rules
% Authors: Antonios Dimakis, John Pavlopoulos, Antonios Anastasopoulos
% Venue:   Findings of ACL 2025
% ===================================================================

\subsection*{dimakis2025dialect: Dialect Normalization using Large Language Models and Morphological Rules}

\noindent \textbf{Authors:} Antonios Dimakis, John Pavlopoulos, Antonios Anastasopoulos \\
\noindent \textbf{Venue:} Findings of ACL 2025 \\[6pt]

Low-resource regional dialects often lack aligned parallel training data required for supervised neural machine translation. In this work, we propose a hybrid dialect normalization framework combining rule-based morphological transformations with few-shot prompting of Large Language Models (LLMs). We evaluate our approach on low-resource Greek dialect proverbs, demonstrating that combining domain rules with LLM reasoning enables effective text standardization without parallel training corpora.

\vspace{12pt}

% ===================================================================
% CITATION KEY: [kresic2024normalizing]
% Title:   Normalizing Swiss German Dialects with the Power of Large Language Models
% Authors: Mihael Kresic, Noorhan Abbas
% Venue:   Procedia Computer Science 2024
% ===================================================================

\subsection*{kresic2024normalizing: Normalizing Swiss German Dialects with the Power of Large Language Models}

\noindent \textbf{Authors:} Mihael Kresic, Noorhan Abbas \\
\noindent \textbf{Venue:} Procedia Computer Science 2024 \\[6pt]

Swiss German encompasses a diverse continuum of regional Alemanic dialects spoken in German-speaking Switzerland. Because Swiss German lacks a standardized written form, digital communication occurs using highly variable phonetic spellings. In this paper, we investigate fine-tuning multilingual pre-trained sequence-to-sequence transformers (mT5-small, mT5-base) on the SwissDial parallel corpus for dialect normalization into Standard High German.

\vspace{12pt}

% ===================================================================
% CITATION KEY: [kuparinen2023dialect]
% Title:   Dialect-to-Standard Normalization: A Large-Scale Multilingual Evaluation
% Authors: Olli Kuparinen, Aleksandra Miletić, Yves Scherrer
% Venue:   Findings of EMNLP 2023
% ===================================================================

\subsection*{kuparinen2023dialect: Dialect-to-Standard Normalization: A Large-Scale Multilingual Evaluation}

\noindent \textbf{Authors:} Olli Kuparinen, Aleksandra Miletić, Yves Scherrer \\
\noindent \textbf{Venue:} Findings of EMNLP 2023 \\[6pt]

Dialect normalization—the task of converting non-standard dialectal text into standard orthography—is a fundamental pre-processing step for low-resource language processing. In this paper, we present a large-scale multilingual benchmark for dialect normalization across four language families covering Finnish, Norwegian, Swiss German, and Slovene dialects. We systematically evaluate character-level SMT, sequence-to-sequence neural architectures, and pre-trained byte-level models (byT5). Our findings demonstrate that sequence-to-sequence neural models consistently outperform character SMT baselines.

\vspace{12pt}

% ===================================================================
% CITATION KEY: [partanen2019dialect]
% Title:   Dialect Text Normalization to Normative Standard Finnish
% Authors: Niko Partanen, Mika Hämäläinen, Khalid Alnajjar
% Venue:   Proceedings of W-NUT @ EMNLP 2019
% ===================================================================

\subsection*{partanen2019dialect: Dialect Text Normalization to Normative Standard Finnish}

\noindent \textbf{Authors:} Niko Partanen, Mika Hämäläinen, Khalid Alnajjar \\
\noindent \textbf{Venue:} Proceedings of W-NUT @ EMNLP 2019 \\[6pt]

We address the problem of normalizing non-standard Finnish dialectal text into Normative Standard Finnish. Regional dialects of Finnish exhibit significant phonological contraction and morphological suffix alterations that cause standard morphological analyzers and POS taggers to fail. We propose a sequence transduction approach operating at character and subword levels across 23 Finnish dialect varieties, demonstrating significant WER reductions.

\vspace{12pt}

% ===================================================================
% CITATION KEY: [gaikwad2025marathi]
% Title:   Marathi Dialect to Pure Marathi using RAG & AI
% Authors: Gaurav Gaikwad, Deepak Panchal, Omprasad Deshmukh, Sonali Patil
% Venue:   IRJET 2025
% ===================================================================

\subsection*{gaikwad2025marathi: Marathi Dialect to Pure Marathi using RAG & AI}

\noindent \textbf{Authors:} Gaurav Gaikwad, Deepak Panchal, Omprasad Deshmukh, Sonali Patil \\
\noindent \textbf{Venue:} IRJET 2025 \\[6pt]

Regional dialects of Marathi, such as Konkani and Malvani, present unique challenges for digital information retrieval systems. In this work, we propose a Retrieval-Augmented Generation (RAG) framework using Sentence Transformers and FAISS vector databases paired with Gemini LLM APIs to normalize Konkani dialect queries into Standard Marathi for automated document question answering.

\vspace{12pt}

% ===================================================================
% CITATION KEY: [pawar2026automatic]
% Title:   Automatic Marathi Dialect Classification in Noisy and Non-Noisy Environments
% Authors: Poonam S. Pawar, Jatinderkumar R. Saini, Shraddha Vaidya
% Venue:   Procedia Computer Science 2026
% ===================================================================

\subsection*{pawar2026automatic: Automatic Marathi Dialect Classification in Noisy and Non-Noisy Environments}

\noindent \textbf{Authors:} Poonam S. Pawar, Jatinderkumar R. Saini, Shraddha Vaidya \\
\noindent \textbf{Venue:} Procedia Computer Science 2026 \\[6pt]

Dialect identification is a critical front-end task for speech and language processing systems operating in multilingual and multi-dialectal regions. We present an acoustic dialect classification system for Marathi sub-dialects (Ahirani, Varhadi, and Standard Marathi) using Mel-Frequency Cepstral Coefficients (MFCC) and Convolutional Neural Networks (CNN).

\vspace{12pt}

% ===================================================================
% CITATION KEY: [mulla2022text]
% Title:   Text-Independent Automatic Dialect Recognition of Marathi Language using Spectro-Temporal Characteristics
% Authors: Rahesha Mulla, B. Suresh Kumar
% Venue:   IJRITCC 2022
% ===================================================================

\subsection*{mulla2022text: Text-Independent Automatic Dialect Recognition of Marathi Language using Spectro-Temporal Characteristics}

\noindent \textbf{Authors:} Rahesha Mulla, B. Suresh Kumar \\
\noindent \textbf{Venue:} IJRITCC 2022 \\[6pt]

Speech features vary significantly across geographical sub-regions of Maharashtra. In this paper, we evaluate text-independent automatic dialect recognition across four Marathi sub-dialects using spectro-temporal voice characteristics (MFCCs, chromagram, spectral centroid) and ANOVA feature selection, achieving 84.24% classification accuracy.

\vspace{12pt}

% ===================================================================
% CITATION KEY: [vyavhare2026marathi]
% Title:   Marathi Dialect Translator
% Authors: V. A. Vyavhare, S. N. Divekar, Harshada Pawar, Karan Kale, Subhash Jagtap
% Venue:   IJSRSET 2026
% ===================================================================

\subsection*{vyavhare2026marathi: Marathi Dialect Translator}

\noindent \textbf{Authors:} V. A. Vyavhare, S. N. Divekar, Harshada Pawar, Karan Kale, Subhash Jagtap \\
\noindent \textbf{Venue:} IJSRSET 2026 \\[6pt]

This study presents a regional Marathi dialect translation web application. Combining rule-based dictionary lookups for common dialectal idioms with neural sequence-to-sequence translation models, the system normalizes regional Marathi dialectal sentences into Standard Marathi.

\vspace{12pt}

% ===================================================================
% CITATION KEY: [amartyaveer2025improving]
% Title:   Improving Dialect Identification in Indian Languages Using Multimodal Features
% Authors: Amartyaveer, Saurabh Kumar, Sumit Sharma, Sathvik Udupa, Sandhya Badiger, Abhayjeet Singh, Deekshitha G, Jesuraja Bandekar, Savitha Murthy, Prasanta Kumar Ghosh
% Venue:   IEEE ICASSP 2025
% ===================================================================

\subsection*{amartyaveer2025improving: Improving Dialect Identification in Indian Languages Using Multimodal Features}

\noindent \textbf{Authors:} Amartyaveer, Saurabh Kumar, Sumit Sharma, Sathvik Udupa, Sandhya Badiger, Abhayjeet Singh, Deekshitha G, Jesuraja Bandekar, Savitha Murthy, Prasanta Kumar Ghosh \\
\noindent \textbf{Venue:} IEEE ICASSP 2025 \\[6pt]

We address the problem of spoken dialect identification across Indian languages using multimodal acoustic and textual representations extracted from dialect-informed ASR systems. Evaluating on the RESPIN dataset, our multimodal fusion approach achieves 79.81% classification accuracy across regional sub-dialects.

\vspace{12pt}

